\chapterbackmatter{Solutions to Weinberg Exercises}
\ExerciseEditorialNote

\begin{enumerate}
\WeinbergSolution{1}
Put
\[
 s=-(p_1+p_2)^2,\qquad
 t=-(p_1-p'_1)^2,\qquad
 u=-(p_1-p'_2)^2,
\qquad s+t+u=4m^2 .
\]
There are three connected second-order contractions, according as the
internal scalar carries \(p_1+p_2\), \(p_1-p'_1\), or \(p_1-p'_2\).
Using Eq.~\eqref{eq:6.2.16} in the position-space result following
Eq.~\eqref{eq:6.1.31}, both vertex integrals can be done:
\[
 \begin{split}
 S^c_{fi}
 ={}&\frac{\ii(2\pi)^{-2}g^2
 \delta^4(p_1+p_2-p'_1-p'_2)}
 {\sqrt{16E_1E_2E'_1E'_2}}\\
 &\cdot\left[
 \frac1{m^2-s-\ii0}
 +\frac1{m^2-t-\ii0}
 +\frac1{m^2-u-\ii0}\right].
 \end{split}
\]
Thus, up to the immaterial overall sign convention for \(\mathcal M\),
\[
 |\mathcal M|^2
 =g^4\left|
 \frac1{m^2-s-\ii0}
 +\frac1{m^2-t-\ii0}
 +\frac1{m^2-u-\ii0}\right|^2 .
\]
For elastic scattering in the center-of-mass frame,
\[
 t=-2k^2(1-\cos\theta),\qquad
 u=-2k^2(1+\cos\theta),\qquad
 k^2=\frac{s}{4}-m^2 .
\]
The two final particles are identical.  If the full solid angle is used,
the differential cross-section therefore includes \(1/2!\):
\[
 \boxed{\;
 \frac{\dd\sigma}{\dd\Omega}
 =\frac{g^4}{128\pi^2s}
 \left|
 \frac1{m^2-s-\ii0}
 +\frac1{m^2+2k^2(1-\cos\theta)}
 +\frac1{m^2+2k^2(1+\cos\theta)}
 \right|^2 .\;}
\]
Equivalently one may omit \(1/2!\) and integrate over only one copy of
the identical-particle phase space.

\WeinbergSolution{2}
It is useful to display the answer with the common fermion numerator
\[
 N(q)=-\ii\gamma_\mu q^\mu+M,\qquad
 D_F(q)=\frac{N(q)}{q^2+M^2-\ii0}.
\]
The factors \(\beta=\ii\gamma^0\) in the propagator convert the external
\(u^\dagger,v^\dagger\) into barred spinors.  Apart from the common
momentum delta function and the scalar external-line factors, the two
connected graphs for antifermion--boson scattering are
\begin{center}
\begin{tikzpicture}[x=0.70cm,y=0.70cm,baseline=-0.5ex,line width=0.65pt,
  ferm/.style={postaction={decorate},decoration={markings,
    mark=at position #1 with {\arrow{Stealth[length=4.6pt,width=3.4pt]}}}},
  ferm/.default=0.55,
  scalar/.style={dash pattern=on 2.8pt off 2.3pt},
  vertex/.style={circle,fill,inner sep=1.25pt},
  every node/.style={font=\small}]
 \coordinate (a) at (-0.95,0); \coordinate (b) at (0.95,0);
 \draw[ferm] (a)--(-2.35,-1.05);
 \draw[ferm] (b)--(a);
 \draw[ferm] (2.35,-1.05)--(b);
 \draw[scalar] (-2.35,1.05)--(a);
 \draw[scalar] (b)--(2.35,1.05);
 \node[vertex] at (a) {}; \node[vertex] at (b) {};
 \node[anchor=north east] at (-2.40,-1.10) {\(\scriptstyle p_1\)};
 \node[anchor=north west] at (2.40,-1.10) {\(\scriptstyle p'_1\)};
 \node[anchor=south east] at (-2.40,1.10) {\(\scriptstyle p_2\)};
 \node[anchor=south west] at (2.40,1.10) {\(\scriptstyle p'_2\)};
\end{tikzpicture}
\qquad
\begin{tikzpicture}[x=0.70cm,y=0.70cm,baseline=-0.5ex,line width=0.65pt,
  ferm/.style={postaction={decorate},decoration={markings,
    mark=at position #1 with {\arrow{Stealth[length=4.6pt,width=3.4pt]}}}},
  ferm/.default=0.55,
  scalar/.style={dash pattern=on 2.8pt off 2.3pt},
  vertex/.style={circle,fill,inner sep=1.25pt},
  every node/.style={font=\small}]
 \coordinate (a) at (-0.95,0); \coordinate (b) at (0.95,0);
 \draw[ferm] (a)--(-2.35,-1.05);
 \draw[ferm] (b)--(a);
 \draw[ferm] (2.35,-1.05)--(b);
 \draw[scalar] (-2.35,1.05)--(a);
 \draw[scalar] (b)--(2.35,1.05);
 \node[vertex] at (a) {}; \node[vertex] at (b) {};
 \node[anchor=north east] at (-2.40,-1.10) {\(\scriptstyle p_1\)};
 \node[anchor=north west] at (2.40,-1.10) {\(\scriptstyle p'_1\)};
 \node[anchor=south east] at (-2.40,1.10) {\(\scriptstyle p'_2\)};
 \node[anchor=south west] at (2.40,1.10) {\(\scriptstyle p_2\)};
\end{tikzpicture}
\end{center}
They give
\[
 \begin{split}
 \mathcal A_{\bar F B}
 ={}&g^2\bar v(p_1)\gamma_5
 \bigl[D_F(-p_1-p_2)+D_F(p'_2-p_1)\bigr]
 \gamma_5v(p'_1).
 \end{split}
\]
The first term is the annihilation-ordering (the crossed version of the
fermion \(s\)-channel graph), and the second is the exchange ordering.
For fermion--antifermion elastic scattering the connected scalar
exchange can be either annihilating or nonannihilating:
\begin{center}
\begin{tikzpicture}[x=0.70cm,y=0.70cm,baseline=-0.5ex,line width=0.65pt,
  ferm/.style={postaction={decorate},decoration={markings,
    mark=at position #1 with {\arrow{Stealth[length=4.6pt,width=3.4pt]}}}},
  ferm/.default=0.55,
  scalar/.style={dash pattern=on 2.8pt off 2.3pt},
  vertex/.style={circle,fill,inner sep=1.25pt},
  every node/.style={font=\small}]
 \coordinate (a) at (0,0.95); \coordinate (b) at (0,-0.95);
 \draw[ferm] (-2.25,0.95)--(a);
 \draw[ferm] (a)--(2.25,0.95);
 \draw[ferm] (b)--(-2.25,-0.95);
 \draw[ferm] (2.25,-0.95)--(b);
 \draw[scalar] (a)--(b);
 \node[vertex] at (a) {}; \node[vertex] at (b) {};
 \node[anchor=south east] at (-2.30,1.00) {\(\scriptstyle p_1\)};
 \node[anchor=south west] at (2.30,1.00) {\(\scriptstyle p'_1\)};
 \node[anchor=north east] at (-2.30,-1.00) {\(\scriptstyle p_2\)};
 \node[anchor=north west] at (2.30,-1.00) {\(\scriptstyle p'_2\)};
\end{tikzpicture}
\qquad
\begin{tikzpicture}[x=0.70cm,y=0.70cm,baseline=-0.5ex,line width=0.65pt,
  ferm/.style={postaction={decorate},decoration={markings,
    mark=at position #1 with {\arrow{Stealth[length=4.6pt,width=3.4pt]}}}},
  ferm/.default=0.55,
  scalar/.style={dash pattern=on 2.8pt off 2.3pt},
  vertex/.style={circle,fill,inner sep=1.25pt},
  every node/.style={font=\small}]
 \coordinate (a) at (-1.05,0); \coordinate (b) at (1.05,0);
 \draw[ferm] (-2.45,1.15)--(a);
 \draw[ferm] (a)--(-2.45,-1.15);
 \draw[scalar] (a)--(b);
 \draw[ferm] (b)--(2.45,1.15);
 \draw[ferm] (2.45,-1.15)--(b);
 \node[vertex] at (a) {}; \node[vertex] at (b) {};
 \node[anchor=south east] at (-2.50,1.20) {\(\scriptstyle p_1\)};
 \node[anchor=north east] at (-2.50,-1.20) {\(\scriptstyle p_2\)};
 \node[anchor=south west] at (2.50,1.20) {\(\scriptstyle p'_1\)};
 \node[anchor=north west] at (2.50,-1.20) {\(\scriptstyle p'_2\)};
\end{tikzpicture}
\end{center}
These give
\[
 \begin{split}
 \mathcal A(F(p_1)+F^c(p_2)\to F(p'_1)+F^c(p'_2))
 ={}&g^2
 \frac{\bar u(p'_1)\gamma_5u(p_1)\,
       \bar v(p_2)\gamma_5v(p'_2)}
 {(p_1-p'_1)^2+m^2-\ii0}\\
 &-g^2
 \frac{\bar v(p_2)\gamma_5u(p_1)\,
       \bar u(p'_1)\gamma_5v(p'_2)}
 {(p_1+p_2)^2+m^2-\ii0}.
 \end{split}
\]
The relative minus sign follows by putting the external fermionic
operators into the same reference order.

Finally, \(F^c+F\to B+B\) has two fermion-exchange graphs:
\begin{center}
\begin{tikzpicture}[x=0.70cm,y=0.70cm,baseline=-0.5ex,line width=0.65pt,
  ferm/.style={postaction={decorate},decoration={markings,
    mark=at position #1 with {\arrow{Stealth[length=4.6pt,width=3.4pt]}}}},
  ferm/.default=0.55,
  scalar/.style={dash pattern=on 2.8pt off 2.3pt},
  vertex/.style={circle,fill,inner sep=1.25pt},
  every node/.style={font=\small}]
 \coordinate (a) at (0,0.82); \coordinate (b) at (0,-0.82);
 \draw[ferm] (-2.05,-1.65)--(b);
 \draw[ferm] (b)--(a);
 \draw[ferm] (a)--(-2.05,1.65);
 \draw[scalar] (a)--(2.05,1.65);
 \draw[scalar] (b)--(2.05,-1.65);
 \node[vertex] at (a) {}; \node[vertex] at (b) {};
 \node[anchor=south east] at (-2.10,1.70) {\(\scriptstyle p_1\)};
 \node[anchor=north east] at (-2.10,-1.70) {\(\scriptstyle p_2\)};
 \node[anchor=south west] at (2.10,1.70) {\(\scriptstyle p'_2\)};
 \node[anchor=north west] at (2.10,-1.70) {\(\scriptstyle p'_1\)};
\end{tikzpicture}
\qquad
\begin{tikzpicture}[x=0.70cm,y=0.70cm,baseline=-0.5ex,line width=0.65pt,
  ferm/.style={postaction={decorate},decoration={markings,
    mark=at position #1 with {\arrow{Stealth[length=4.6pt,width=3.4pt]}}}},
  ferm/.default=0.55,
  scalar/.style={dash pattern=on 2.8pt off 2.3pt},
  vertex/.style={circle,fill,inner sep=1.25pt},
  every node/.style={font=\small}]
 \coordinate (a) at (0,0.82); \coordinate (b) at (0,-0.82);
 \draw[ferm] (-2.05,-1.65)--(b);
 \draw[ferm] (b)--(a);
 \draw[ferm] (a)--(-2.05,1.65);
 \draw[scalar] (a)--(2.05,1.65);
 \draw[scalar] (b)--(2.05,-1.65);
 \node[vertex] at (a) {}; \node[vertex] at (b) {};
 \node[anchor=south east] at (-2.10,1.70) {\(\scriptstyle p_1\)};
 \node[anchor=north east] at (-2.10,-1.70) {\(\scriptstyle p_2\)};
 \node[anchor=south west] at (2.10,1.70) {\(\scriptstyle p'_1\)};
 \node[anchor=north west] at (2.10,-1.70) {\(\scriptstyle p'_2\)};
\end{tikzpicture}
\end{center}
They give
\[
 \begin{split}
 \mathcal A_{\bar F F\to BB}
 =g^2\bar v(p_1)\gamma_5\bigl[
 D_F(p_2-p'_1)+D_F(p_2-p'_2)\bigr]\gamma_5u(p_2).
 \end{split}
\]
With the chapter's delta-normalized states, the connected matrix element
for any of these processes is
\[
 S^c_{fi}=-\ii(2\pi)^{-2}\delta^4(P_f-P_i)
 \prod_{\text{external }B}(2E_B)^{-1/2}\,\mathcal A ,
\]
where the displayed spinors already contain the fermionic external-line
normalizations.  These formulas contain no remaining position or internal
momentum integrals.

\WeinbergSolution{3}
At first order, all \(4!\) attachments of the four external scalars
cancel the \(1/4!\) in \(H_I\).  Hence
\[
 S^{(1)c}_{fi}
 =\frac{-\ii(2\pi)^{-2}g\,
 \delta^4(p_1+p_2-p'_1-p'_2)}
 {\sqrt{16E_1E_2E'_1E'_2}},
 \qquad \mathcal M^{(1)}=-g .
\]
The center-of-mass differential cross-section is consequently
\[
 \boxed{\quad
 \frac{\dd\sigma}{\dd\Omega}
 =\frac{g^2}{128\pi^2s}\quad}
\]
when the full solid angle is retained for the identical final particles.

At order \(g^2\), vacuum bubbles and disconnected spectator pieces do
not contribute to the connected \(2\to2\) matrix element.  The three
remaining bubble diagrams have identical symmetry factor \(1/2\).
Defining
\[
 I(P)=\int\frac{\dd^4q}{(2\pi)^4}\,
 \frac1{(q^2+m^2-\ii0)((q+P)^2+m^2-\ii0)},
\]
the correction can be written, with all \(x\)-integrals done, as
\[
 \mathcal M^{(2)}
 =\frac{\ii g^2}{2}\left[
 I(p_1+p_2)+I(p_1-p'_1)+I(p_1-p'_2)\right],
\]
up to the same overall \(\ii\mathcal M\) convention used at first order.
Equivalently,
\[
 S^{(2)c}_{fi}
 =\frac{-\ii(2\pi)^{-2}\delta^4(P_f-P_i)}
 {\sqrt{16E_1E_2E'_1E'_2}}\,
 \mathcal M^{(2)} .
\]
Each term contains precisely one independent four-momentum, in agreement
with \(L=I-V+1=2-2+1=1\).  The integrals are ultraviolet divergent; the
problem asks only for their unrenormalized one-momentum representation.

\WeinbergSolution{4}
Differentiation commutes with forming the contraction, provided the
derivative acts on the \(x\) argument:
\[
 \bigl(\partial_\mu\psi_\ell(x)\,\bar\psi_m(y)\bigr)_0
 =\partial_{x^\mu}\bigl(\psi_\ell(x)\bar\psi_m(y)\bigr)_0
 =-\ii\,\partial_{x^\mu}\Delta_{\ell m}(x-y).
\]
For a Dirac field,
\[
 \Delta_{\ell m}(x-y)
 =\int\frac{\dd^4q}{(2\pi)^4}\,
 \frac{\bigl[(-\ii\gamma_\nu q^\nu+M)\beta\bigr]_{\ell m}
 e^{\ii q\cdot(x-y)}}{q^2+M^2-\ii0}.
\]
Therefore the line with the differentiated end contributes
\[
 -\ii\int\frac{\dd^4q}{(2\pi)^4}\,
 \frac{\ii q_\mu
 \bigl[(-\ii\gamma_\nu q^\nu+M)\beta\bigr]_{\ell m}}
 {q^2+M^2-\ii0}\,e^{\ii q\cdot(x-y)}.
\]
Thus the ordinary Dirac propagator factor is simply multiplied by
\(\ii q_\mu\) when the momentum \(q\) flows from \(y\) to \(x\).  No extra
contact term occurs for one derivative; contact terms first arise when
enough time derivatives act on the step functions.

\WeinbergSolution{5}
Let \(\phi\) denote the free interaction-picture field and write
\[
 D(x-y)\equiv\bra{\Omega}T\{\phi(x)\phi(y)\}\ket{\Omega}
 =-\ii\Delta_F(x-y).
\]
The theorem of Section 6.4, with vacuum-bubble normalization understood,
gives the Gell-Mann--Low ratio
\[
 \bra{\Omega}T\{\Phi(x_1)\cdots\Phi(x_n)\}\ket{\Omega}
 =
 \frac{\bra{\Omega}T\{\phi(x_1)\cdots\phi(x_n)
 e^{-\ii(g/3!)\int \dd^4z\,\phi(z)^3}\}\ket{\Omega}}
 {\bra{\Omega}T\{e^{-\ii(g/3!)\int \dd^4z\,\phi(z)^3}\}\ket{\Omega}} .
\]
At order \(g\), three equivalent contractions attach \(x\) to the
vertex and contract the remaining two fields together.  Hence
\[
 \boxed{\quad
 \bra{\Omega}\Phi(x)\ket{\Omega}
 =-\frac{ig}{2}\int \dd^4z\,D(x-z)D(0)+O(g^3).\quad}
\]
The absence of an order-\(g^2\) term follows from the odd total number of
free fields.

For the time-ordered two-point function, the denominator cancels every
vacuum bubble.  Through order \(g^2\) the remaining Wick contractions are
\[
 \begin{split}
 \bra{\Omega}T\{\Phi(x)\Phi(y)\}\ket{\Omega}
 ={}&D(x-y)\\
 &+\frac{(-ig)^2}{2}\int \dd^4z\,\dd^4w\,
 D(x-z)D(z-w)^2D(w-y)\\
 &+\frac{(-ig)^2}{2}\int \dd^4z\,\dd^4w\,
 D(x-z)D(y-z)D(z-w)D(0)\\
 &+\frac{(-ig)^2}{4}\int \dd^4z\,\dd^4w\,
 D(x-z)D(0)D(y-w)D(0)+O(g^4).
 \end{split}
\]
The second line is the one-loop self-energy graph, the third is the
connected tadpole insertion, and the fourth is the disconnected product
\(\bra{\Omega}\Phi(x)\ket{\Omega}
 \bra{\Omega}\Phi(y)\ket{\Omega}\) to this order.  The coefficients
\(1/2,1/2,1/4\) are the corresponding Wick-contraction symmetry factors.
\end{enumerate}
